\documentclass[10pt,a4paper]{amsart}
\usepackage[margin=19mm]{geometry}
\usepackage{amsmath,amssymb,mathtools}
\usepackage[scr=rsfs,cal=euler]{mathalfa}
\usepackage{xfrac}
\usepackage[unicode,hidelinks,bookmarks=false]{hyperref}
\newcommand{\ie}{i.e.\ }
\newcommand{\cf}{cf.\ }
\newcommand{\resp}{respectively\ }
\newcommand{\Subsection}{\subsection}
\providecommand{\nobreakdash}{\nobreak\hbox{-}}
\setlength{\emergencystretch}{2em}
\multlinegap0pt
\usepackage{bm}
\usepackage{bbm}
\usepackage{dsfont}
\usepackage{xspace}
\usepackage{cancel}
\usepackage{mathtools}
\usepackage{amsthm}
\makeatletter
\@ifundefined{lemma}{\newtheorem{lemma}{Lemma}}{}
\@ifundefined{proposition}{\newtheorem{proposition}{Proposition}}{}
\@ifundefined{theorem}{\newtheorem{theorem}{Theorem}}{}
\makeatother
\DeclareMathOperator{\gl}{GL}
\DeclareMathOperator{\spn}{Span}
\newcommand{\F}{\mathbb{F}}
\newcommand{\Ho}{H_\circ}
\newcommand{\Uo}{U_\circ}
\newcommand{\Wo}{W_\circ}
\newcommand{\defeq}{\vcentcolon=}
\newcommand{\eps}{\epsilon}
\newcommand{\ho}{H_\circ}
\newcommand{\lb}{\lambda}
\newcommand{\subs}{\subset}
\newcommand{\uo}{U_\circ}
\newcommand{\wo}{W_\circ}
\renewcommand{\b}{\textbf{b}}
\title[Optimal s-boxes against alternative operations and linear pr]{Optimal s-boxes against alternative operations and linear propagation}
\author{Marco Calderini, Roberto Civino, Riccardo Invernizzi}
\date{Source: arXiv:2403.20059v3 (CC BY 4.0)}
\begin{document}
\maketitle

\noindent\emph{Source and licence.} Optimal s-boxes against alternative operations and linear propagation. Version arXiv:2403.20059v3 (CC BY 4.0), distributed under the Creative Commons Attribution 4.0 International licence, as recorded in the arXiv licence metadata for this exact version and shown on the abstract page https://arxiv.org/abs/2403.20059v3. Full rights notice: licenses/arxiv-2403.20059-CC-BY-4.0.txt.\par\smallskip

\noindent\textit{Editorial scope.} Counted unit: the Theorem whose source label is \texttt{thm:caratt\_h0\_msomme} in the article (source0.tex lines 832--889, source label \texttt{thm:caratt\_h0\_msomme}), delivered together with the article's own complete original proof of it (source0.tex lines 890--943, one \texttt{proof} environment of 54 lines). No mathematics is written, repaired or re-derived: statement and proof are the article's own text line for line. Required context retained uncounted: prop:info\_prod (lines 679--693); lem:spazi\_fissi (lines 719--722). Every definition, display and label of those lines is retained unchanged. Omitted from this package and not redistributed: 1--678, 694--718, 723--831, 944--1425. Nothing inside a retained excerpt is omitted or altered: every retained body is delivered exactly as the article's lines are written, so no omission receipt of any kind accompanies this package. Citations: the retained excerpts carry no citation command at all, so no bibliography entry of the article is transcribed and no reference is redistributed. Reference adaptation: none was needed and none was made; every reference inside the delivered unit resolves inside it, so no unresolved reference or citation is produced. Printed slips kept literal: no printed word, symbol, hypothesis or number of the counted statement or of its proof was altered. Layout and class: the article's own class is \texttt{llncs} with packages including amsmath, upquote, orcidlink, babel, regexpatch, cfr-lm, fontenc, microtype, graphicx, diagbox, xcolor, listings, csquotes, booktabs; the wrapper is the lane's pinned \texttt{amsart} editorial wrapper at 10pt on A4, which restates the theorem-like environments the retained text needs and adds the article's own macro definitions that the retained text needs (\texttt{\textbackslash F}, \texttt{\textbackslash Ho}, \texttt{\textbackslash Uo}, \texttt{\textbackslash Wo}, \texttt{\textbackslash b}, \texttt{\textbackslash defeq}, \texttt{\textbackslash eps}, \texttt{\textbackslash gl}, \texttt{\textbackslash ho}, \texttt{\textbackslash lb}, \texttt{\textbackslash spn}, \texttt{\textbackslash subs}, \texttt{\textbackslash uo}, \texttt{\textbackslash wo}), with extra wrapper packages (bm, bbm, dsfont, xspace, cancel, mathtools). The delivered numbering is the unit's own. Assets: no retained span contains a figure environment, a graphics-inclusion command, a PDF-page inclusion, a table or a TikZ picture; no image file is copied into this package, the package declares no asset dependency and no image file is redistributed with it. Licence: version arXiv:2403.20059v3 of this article is distributed under the Creative Commons Attribution 4.0 International licence, as recorded in the arXiv licence metadata for this exact version and shown on the abstract page https://arxiv.org/abs/2403.20059v3. A full rights notice is delivered at licenses/arxiv-2403.20059-CC-BY-4.0.txt. Characters: every retained excerpt is pure ASCII, so the delivered engine, XeLaTeX with its default Latin Modern text font, reports no missing character.
\begin{proposition}
	\label{prop:info_prod}
    Let $\circ$ be an operation as above and let $\cdot$ be the induced product. Then
    \begin{itemize}
        \item $\cdot$ is distributive w.r.t\ $+$, i.e.\ $(V, +, \cdot)$ is an $\F_2$-algebra;
        \item $\Uo \subs \wo$;
        \item for each $x,y \in V$ there exists $\eps_{x,y} \in \uo$ such that
        \[
        x + y = x \circ y + \eps_{x, y},
        \]
        with $\eps_{x, y} = x \cdot y = (0, \dots, 0, (x_1,\dots, x_{n-d})E_y)$;
        \item $\uo$ is composed of all possible vectors $w \in \wo$ whose last $d$ components are $\F_2$-linear combinations of the defining vectors $\b_{ij}$;
        \item $x \cdot y \cdot z = 0$, for all $x, y, z \in V$.
    \end{itemize}
\end{proposition}

\begin{lemma}
	\label{lem:spazi_fissi}
	For each $\lb \in \Ho$ it holds $\Wo\lb = \Wo$ and $\Uo\lb = \Uo$.
\end{lemma}

\begin{theorem}
	\label{thm:caratt_h0_msomme}
	Let $\circ =(\circ_{1}, \circ_{2}, \dots, \circ_{b})$ be a parallel alternative operation on $V$ such that for each $1\leq j \leq b$ $\circ_j$ is an alternative operation on $V_j$. Let us assume that every $\circ_{j}$ is such that $\dim(W_{\circ_j}) = s-2$ and it is defined by a vector $\b \in (\F_2)^{s-2}$.
	Let $\lb \in (\F_2)^{n\times n}$. Then, $\lb \in H_\circ$ if and only if it can be represented in the block form
	\[ 
	\lambda = 
	\left(\begin{array}{@{}c|c|c@{}}
		\begin{matrix}
			A_{11} & B_{11} \\
			C_{11} & D_{11}
		\end{matrix}
		&
		\cdots
		&
		\begin{matrix}
			A_{1b} & B_{1b} \\
			C_{1b} & D_{1b}
		\end{matrix}
		\\
		\hline
		\vdots
		&
		\ddots
		&
		\vdots
		\\
		\hline
		\begin{matrix}
			A_{b1} & B_{b1} \\
			C_{b1} & D_{b1}
		\end{matrix}
		&
		\cdots
		&
		\begin{matrix}
			A_{bb} & B_{bb} \\
			C_{bb} & D_{bb}
		\end{matrix}
		\\
	\end{array}\right),
	\]
	where
	\begin{enumerate}
		\item $A_{ij} \in (\F_2)^{2\times 2}$ such that for each row and each column of blocks there exists one and only one nonzero $A_{ij}$; moreover, all the nonzero $A_{ij}$ are invertible;
		\item $B_{ij} \in (\F_2)^{2 \times (s-2)}$;
		\item $C_{ij} = \mathbb{0}_{(s-2) \times 2}$;
		\item $D_{ij} \in (\F_2)^{(s-2) \times (s-2)}$ such that if $A_{ij}$ is zero, then $\b D_{ij}=\textbf{0}$, and if $A_{ij}$ is invertible, then $\b D_{ij} = \b$. Moreover, the matrix $D$ defined by
		\[
		D \defeq
		\begin{pmatrix}
			D_{11} & \cdots & D_{1b} \\
			\vdots & \ddots & \vdots \\
			D_{b1} & \cdots & D_{bb}
		\end{pmatrix}
		\]
		is invertible.
	\end{enumerate}
\end{theorem}

\begin{proof}
%Let us first introduce some notation in order to simplify the proof. We denote with $e_i^j$ the vector with a one in the $i$-th component of the $j$-th block, i.e. what would be $e_{s(j-1) + i}$. In this way, the $i$-th strong component (the component spanned by strong keys in the $i$-th block is $\spn\{ e^i_1, e^i_2 \}$ and the $i$-th weak component (also denoted by $\wo^i$) is $\spn\{ e^i_3, ..., e^i_s \}$. We denote $\wo = \wo^1 \oplus ...\oplus \wo^b$.
%    
%$(\Rightarrow)$ Let $\lb \in \ho$. We have to prove that if we divide it in blocks as in the statements, all the blocks will have the desired properties. Notice that by definition the block $A_{ij}$ sends the $i$-th strong component into the $j$-th strong component, the block $B_{ij}$ sends the $i$-th strong component into the $j$-th weak component, the block $C_{ij}$ sends the $i$-th weak component into the $j$-th strong component while the block $D_{ij}$ sends the $i$-th weak component into the $j$-th weak component. 
%
%First of all, thanks to Lemma \ref{lem:spazi_fissi} (which can be naturally extended to the parallel setting) we have $\wo \lb = \wo$. In light of what we just observed, this implies $C_{ij} = 0$ for all $i,j$. As a consequence, $\wo$ is contained in the image of $D$; since $\dim(\wo) = b(s-2)$, $D$ must be invertible.
%
%Let us now focus on the $A_{ij}$. 
%If we take a vector $x_i$ whose nonzero strong components are contained in $V_i$ we want to prove that $x_i\lambda = x_j$ where $x_j$ is a vector whose nonzero strong components are contained in $V_j$ for some $j$.
%We know that $e^i_1 \cdot e^i_2 \neq 0$; on the other hand, $e^j_{i_1} \cdot e^k_{i_2} = 0$ if $j \neq k$, for all combinations of $i_1, i_2 \in \{1, 2\}$. 
%Let us look at $\spn \{e^i_1, e^i_2 \} \lb$, and assume that it contains an element with more than one nonzero strong component. Assume for instance $e_1^i \lb = e^j_{r} + e^k_s$. Since we started from a subspace of dimension 2, its image cannot span both the strong subspaces of $V_{j}$ and $V_{k}$. Let us assume that it does not generate $V_k$. Then we will need another basis vector $e^l_t$ such that $l \neq i$ and $e_l^t \lb$ has nonzero strong component in $V_k$. But this implies
%\[
%0 = e_1^i \cdot e_l^t = (e_1^i \cdot e_l^t) \lb = e_1^i\lb \cdot e_l^t \lb \neq 0,
%\]
%which is absurd. Hence, each vector with only one nonzero strong component is sent into a vector with only one nonzero strong component, i.e. for each row (or column) there is exactly one nonzero $A_{ij}$. Moreover, the nonzero $A_{ij}$ must span the $j$-th strong subspace, and must then be invertible.
%
%Now recall that since $d = s-2$, $e^i_1 \cdot e^i_2$ is the vector with $(0, 0, \b)$ in $V_i$ and $0$ elsewhere. Hence, the component in $V_j$ of $e^i_1 \cdot e^i_2$ is $\b D_{ij}$. If $A_{ij} = 0$,it must be equal to $e^i_1 A_{ij} \cdot e^i_2 A_{ij} = 0$, which implies $\b D_{ij} = 0$. On the other hand, if $A_{ij}$ is invertible we must have $\b D_{ij} = \b$, due to Lemma \ref{lem:spazi_fissi} again.	  
% 
%$(\Leftarrow)$ We now want to show that if $\lb$ is as in the statement, $\lb \in \ho = \hom(V, +, \cdot)$. Since by construction $\lb \in \gl_+(V)$, we only have to check that $(x \cdot y)\lb = x\lb \cdot y\lb$ for all $x, y \in V$. Moreover, since by Proposition \ref{prop:info_prod} the dot product is distributive with respect to $+$, we may restrict to check the equality component-wise. 
% 
%If at least one among $x$ and $y$ is a weak key, both sides of the equation are zero. Indeed, assume $x$ is a weak key. On the left hand side $x \cdot y = 0$, by definition of dot product. On the right hand side, since all the $C_{ij} = 0$, $x\lb \in \wo$. Hence again, $x\lb \cdot y\lb = 0$. 
% 
%There are hence only two cases left to consider: strong vectors coming from different blocks and strong vectors coming from the same block. We start from the first case, namely $(e^j_{i_1} \cdot e^k_{i_2})\lb = e^j_{i_1} \lb \cdot e^k_{i_2}\lb$ for $j \neq k$ and $i_1, i_2 \in \{1, 2\}$.
%As already observed, $j \neq k$ implies $e^j_{i_1} \cdot e^k_{i_2} = 0$. On the right hand side, the condition on the $A$ blocks of $\lb$ ensures that the part of $e^j_{i_1} \lb$ and $e^k_{i_2}\lb$ that lies outside $\wo$ end up in two different blocks, namely the blocks $V_l$ such that $A_{jl} \neq 0$ and $A_{kl} \neq 0$ respectively. This proves $e^j_{i_1} \lb \cdot e^k_{i_2}\lb = 0$.
%Notice that here we do not have any restriction on the components of $e^j_{i_1} \lb$ and $e^k_{i_2}\lb$ in $\wo$, since we impose no restriction on the $B$ blocks of $\lb$. However by definition these components never contribute to the dot product.
%
%Now look at the second case: strong vectors coming from the same block, i.e. $(e^j_{i_1} \cdot e^j_{i_2})\lb = e^j_{i_1} \lb \cdot e^j_{i_2}\lb$. If $i_1 = i_2$ both sides of the equality are $0$, so we may assume without loss of generality $i_1 = 1$ and $i_2 = 2$. Then by definition $e_1^j \cdot e_2^j$ is a vector with $(0, 0, \b)$ in the components corresponding to $V_j$, and zero elsewhere. For any given $j$, there will be only one index $k$ such that $\b D_{jk} = \b$, while for all $l \neq k$ we have $\b D_{jl} = 0$. Hence the left hand side is a vector with $(0, 0, \b)$ in the components of $V_k$ and zero elsewhere. But then the same $k$ is such that $A_{jk}$ is invertible, and $A_{jl} = 0$ for all $l \neq k$. Hence $e_1^j\lb$ and $e_2^j\lb$ are two distinct vectors in the subspace $V_k$, and their dot product is a vector with $(0, 0, \b)$ in the same component, i.e. the equality holds. This completes the proof.

First, let us introduce some notation to simplify the proof. We denote \( e_i^j \) as the vector with a one in the \( i \)-th component of the \( j \)-th block, which is equivalent to \( e_{s(j-1) + i} \). Thus, the \( i \)-th strong component (the component spanned by nonweak keys in the \( i \)-th block) is \( \spn\{ e^i_1, e^i_2 \} \), and the \( i \)-th weak component (also denoted as \( \wo^i \)) is \( \spn\{ e^i_3, \dots, e^i_s \} \). We define \( \wo = \wo^1 \oplus \cdots \oplus \wo^b \).\\

\noindent (\(\Rightarrow\)) Let \( \lambda \in \ho \). We need to show that if \( \lambda \) is divided into blocks as described, all the blocks will have the desired properties. Notice that by definition, the block \( A_{ij} \) maps the \( i \)-th strong component to the \( j \)-th strong component, the block \( B_{ij} \) maps the \( i \)-th strong component to the \( j \)-th weak component, the block \( C_{ij} \) maps the \( i \)-th weak component to the \( j \)-th strong component, and the block \( D_{ij} \) maps the \( i \)-th weak component to the \( j \)-th weak component.

Firstly, by Lemma \ref{lem:spazi_fissi} (which extends naturally to the parallel setting), we have \( \wo \lambda = \wo \). This implies that \( C_{ij} = 0 \) for all \( i, j \). Consequently, \( \wo \) is contained in the image of \( D \); since \( \dim(\wo) = b(s-2) \), \( D \) must be invertible.

Now, let us focus on the blocks \( A_{ij} \). If we take a vector \( x_i \) with nonzero strong components in \( V_i \), we want to prove that \( x_i \lambda = x_j \) where \( x_j \) is a vector with nonzero strong components in \( V_j \) for some \( j \). We know that \( e^i_1 \cdot e^i_2 \neq 0 \); however, \( e^j_{i_1} \cdot e^k_{i_2} = 0 \) if \( j \neq k \), for all combinations of \( i_1, i_2 \in \{1, 2\} \).

Consider \( \spn \{e^i_1, e^i_2 \} \lambda \), and assume it contains an element with more than one nonzero strong component. Without loss of generality, suppose that the strong component of \( e_1^i \lambda\) is given by \(e^j_{r} + e^k_s \),  for some $k\ne j$ and \(r,s\in\{1,2\}\). Since we started with a 2-dimensional subspace, its image cannot span both strong subspaces of \( V_j \) and \( V_k \). If it does not generate \( V_k \), we would need another basis vector \( e^l_t \) with \( l \neq i \) such that \( e^l_t \lambda \) has a nonzero strong component in \( V_k \) different from \(e^k_s \). This implies
\[
0 = e_1^i \cdot e^l_t = (e_1^i \cdot e^l_t) \lambda = e_1^i \lambda \cdot e^l_t \lambda \neq 0,
\]
which is a contradiction. Therefore, each vector with only one nonzero strong component is mapped to a vector with only one nonzero strong component, meaning that for each row (or column) there is exactly one nonzero \( A_{ij} \). Furthermore, the nonzero \( A_{ij} \) must span the \( j \)-th strong subspace and must be invertible.

Since \( d = s-2 \), \( e^i_1 \cdot e^i_2 \) is the vector with \((0, 0, \b)\) in \( V_i \) and zero elsewhere. Therefore, the component in \( V_j \) of \( e^i_1 \cdot e^i_2 \) is \( \b D_{ij} \). If \( A_{ij} = 0 \), it must be that \( e^i_1 A_{ij} \cdot e^i_2 A_{ij} = 0 \), which implies \( \b D_{ij} = 0 \). Conversely, if \( A_{ij} \) is invertible, then \( \b D_{ij} = \b \) due to Lemma \ref{lem:spazi_fissi}.\\


\noindent(\(\Leftarrow\)) We now need to show that if \( \lambda \) has the form described in the statement, then \( \lambda \in \ho\). Since \( \lambda \in \gl(V,+) \) by construction, we only need to verify that \( (x \cdot y)\lambda = x\lambda \cdot y\lambda \) for all \( x, y \in V \). Moreover, since by Proposition \ref{prop:info_prod} the dot product is distributive with respect to \( + \), we can check this equality component-wise.

If either \( x \) or \( y \) is a weak key, both sides of the equation are zero. If \( x \) is a weak key, then \( x \cdot y = 0 \) by definition of the dot product. On the right-hand side, since all \( C_{ij} = 0 \), \( x\lambda \in \wo \), and thus \( x\lambda \cdot y\lambda = 0 \).

We are left with two cases: strong vectors from different blocks and strong vectors from the same block. For the first case, consider \( (e^j_{i_1} \cdot e^k_{i_2})\lambda = e^j_{i_1} \lambda \cdot e^k_{i_2} \lambda \) for \( j \neq k \) and \( i_1, i_2 \in \{1, 2\} \). Since \( j \neq k \) implies \( e^j_{i_1} \cdot e^k_{i_2} = 0 \), the condition on the \( A \) blocks of \( \lambda \) ensures that the part of \( e^j_{i_1} \lambda \) and \( e^k_{i_2}\lambda \) that lies outside \( \wo \) ends up in different blocks. Thus, \( e^j_{i_1} \lambda \cdot e^k_{i_2}\lambda = 0 \). No restriction is imposed on the components of \( e^j_{i_1} \lambda \) and \( e^k_{i_2}\lambda \) in \( \wo \), as the \( B \) blocks of \( \lambda \) are unrestricted. However, these components do not affect the dot product.

For the second case, strong vectors from the same block, consider \( (e^j_{i_1} \cdot e^j_{i_2})\lambda = e^j_{i_1} \lambda \cdot e^j_{i_2} \lambda \). If \( i_1 = i_2 \), both sides are zero. Without loss of generality, assume \( i_1 = 1 \) and \( i_2 = 2 \). Then, \( e_1^j \cdot e_2^j \) is a vector with \( (0, 0, \b) \) in the components corresponding to \( V_j \), and zero elsewhere. For a given \( j \), there is exactly one index \( k \) such that \( \b D_{jk} = \b \), and for all \( l \neq k \), \( \b D_{jl} = 0 \). Thus, the left-hand side is a vector with \( (0, 0, \b) \) in the components of \( V_k \) and zero elsewhere. Since \( A_{jk} \) is invertible and \( A_{jl} = 0 \) for \( l \neq k \), \( e_1^j \lambda \) and \( e_2^j \lambda \) are distinct vectors in the subspace \( V_k \), and their dot product is a vector with \( (0, 0, \b) \) in the same component, proving the equality. This completes the proof.\qed

\end{proof}

\end{document}
